\section{Related Work}
\label{sec:related}

\subsection{LLM-Based Recommendation}

The integration of large language models into recommendation systems has
followed several paradigms. \textbf{Embedding-based methods} such as
LLMEmb~\cite{llmemb} extract item representations from LLM hidden states and
compute similarity scores. \textbf{Sequential methods} such as
LLM2Rec~\cite{llm2rec} and LLM-ESR~\cite{llmesr} combine LLM features with
sequential recommenders such as SASRec~\cite{sasrec}. \textbf{Profile-based
methods} including ProEx~\cite{proex} and ProMax~\cite{promax} use LLM-derived
profiles as preference signals. \textbf{High-order, representation, and intent
methods} such as ELMRec~\cite{elmrec}, RLMRec~\cite{rlmrec}, and
IRLLRec~\cite{irllrec} combine collaborative structure or user intent with
language-derived representations. A complementary line treats the LLM itself as
a zero-shot ranker over a small candidate set, formatting the user history and
candidates in a prompt and reading a ranked list back out~\cite{llmrank}.

Because these methods differ in architecture and training objective, fair
comparison requires a tightly controlled protocol. We treat them as official
external baselines: each method keeps its declared architecture and objective,
while the data adapter, LLM backbone, and score exporter are standardized so
all methods are evaluated over the same candidate rows.
We differ by holding the candidate set, backbone, and importer fixed across all
eight baselines and our method, so the comparison isolates the scoring signal
rather than the surrounding architecture.

\subsection{LLMs as Scorers and Rerankers}

A parallel line of work in information retrieval uses LLMs directly as
relevance scorers or list rerankers rather than as feature extractors.
Listwise approaches prompt the model to reorder a candidate list in a single
pass, either through a sliding-window permutation generation strategy such as
RankGPT~\cite{rankgpt} or by reading a full listwise ranking out of one
prompt as in LRL~\cite{lrl}. Pairwise prompting (PRP) instead asks the model
which of two candidates is more relevant and aggregates the comparisons into a
ranking, trading inference cost for robustness~\cite{prp}. A related strand
studies whether LLMs can stand in for human relevance assessors at all, framing
the model as a (potentially noisy) judge of query--document
relevance~\cite{faggioli2023perspectives}. These methods establish that LLM
relevance signals are competitive, but they are predominantly listwise or
pairwise and are evaluated on document-retrieval benchmarks rather than
sequential recommendation.
We differ by using a task-grounded \emph{pointwise} per-candidate relevance
posterior in the recommendation setting, scoring each candidate independently
under a fixed same-candidate contract instead of a listwise or pairwise prompt.

\subsection{Calibration and Verbalized Confidence}

Calibration of neural predictions has been extensively studied in supervised
learning~\cite{guo2017calibration,niculescu2005predicting}. Post-hoc methods
such as temperature scaling, Platt scaling~\cite{platt1999probabilistic}, and
isotonic regression~\cite{zadrozny2002transforming} can improve probability
calibration when a suitable calibration set is available. In the LLM setting,
prior work has studied whether models can express their own uncertainty in
words or probabilities~\cite{kadavath2022language,tian2023just,lin2022teaching},
and how best to \emph{elicit} that confidence: black-box prompting strategies
can surface a model's verbalized uncertainty even without access to its
logits~\cite{xiong2024can}. This literature largely asks whether a scalar
confidence is well calibrated against correctness, typically in classification
or question answering.
We differ by asking not whether a confidence value is calibrated in isolation
but whether an elicited, task-grounded uncertainty signal changes a downstream
ranking decision under a fixed candidate set.

\subsection{Uncertainty, Conformal, and Selective Prediction}

Uncertainty quantification has been studied across ranking, recommendation, and
prediction-set frameworks. For neural learning-to-rank, prior work measures the
calibration and uncertainty of ranking models directly~\cite{penha2021calibration},
while field-aware calibration targets calibrated probability estimates in
large-scale prediction settings~\cite{pan2020fieldaware}. A separate and
frequently confused notion is \emph{calibrated recommendations} in the sense of
Steck~\cite{steck2018calibrated}, which matches the genre/category
\emph{distribution} of a recommendation list to the user's consumption history;
this is distribution calibration of list composition, a different sense from the
probability calibration of a per-candidate score that we study, and we cite it
only to disambiguate the term. Conformal and selective-prediction methods
provide distribution-free guarantees: conformal prediction yields valid
prediction sets~\cite{angelopoulos2021gentle}, conformal risk control extends
these guarantees to bounded monotone losses~\cite{angelopoulos2024crc}, and
selective classification studies when a predictor should abstain to control
risk, whether through classical foundations~\cite{elyaniv2010foundations},
deep selective classifiers~\cite{geifman2017selective}, or an explicit
risk--coverage training objective such as SelectiveNet~\cite{geifman2019selectivenet}.
We differ by requiring every method to rank every candidate -- there is no
abstention and no prediction set -- so we use these formulations only to
motivate decision-useful uncertainty, and we report that under our protocol the
risk-adjusted family does not improve ranking.

\subsection{Sampled versus Full-Catalog Evaluation}

Offline recommendation metrics are sensitive to how the candidate set is
constructed. Sampled candidate evaluation can distort metric estimates and
their relative ordering~\cite{krichene2020sampled}, and the specific choice of
which target items and negatives to sample further shifts measured
performance~\cite{canamares2020target}. These findings caution against mixing
full-catalog retrieval metrics with reranking metrics computed over different
candidate pools, or comparing methods that score different candidates.
We differ by enforcing a same-candidate score contract -- every method emits a
finite score for the exact candidate key set and all metrics are computed by a
shared importer -- which makes the reranking claim auditable while explicitly
scoping it away from full-catalog SOTA.
